% =============================================================================
% 6.7 SEGURIDAD DEL MOTOR Y POLÍTICAS RLS
% =============================================================================
\section{Seguridad del Motor y Políticas RLS}
\label{sec:rls}

La seguridad de los datos no descansa únicamente en el backend: PostgreSQL aplica
\textbf{Row Level Security} (RLS) directamente en el motor, de modo que incluso un acceso
directo desde el cliente web (con la \texttt{anon key} de Supabase) queda restringido a
las filas que el usuario tiene derecho a ver. Esta es una defensa en profundidad: aunque
fallara una comprobación en la aplicación, el motor seguiría protegiendo los datos.

\subsection{Configuración de Políticas de Seguridad a Nivel de Fila (RLS)}
\label{subsec:politicas_rls}

De las 25 tablas de \texttt{core}, \textbf{23 tienen RLS habilitado}. Las dos excepciones
(\comando{portfolio\_items} y \\ \comando{talent\_score\_snapshots}) se acceden
exclusivamente a través del backend o de funciones, nunca directamente desde el cliente.
La figura~\ref{fig:rls-flujo} ilustra los dos caminos de acceso a los datos y cómo RLS
gobierna el acceso directo del cliente.

\clearpage
\vspace{0.3cm}
\begin{figure}[h!]
\centering
\begin{tikzpicture}[node distance=0.9cm and 1.6cm,
    box/.style={rectangle, rounded corners=3pt, draw, font=\sffamily\scriptsize, align=center,
        minimum width=2.8cm, minimum height=0.9cm, inner sep=3pt},
    rel/.style={-{Stealth[length=2mm]}, semithick}]
    \node[box, fill=c4Componente!40] (cli) at (0,2.5) {Cliente SPA\\\scriptsize (anon key)};
    \node[box, fill=c4Componente!40] (be) at (0,0) {Backend\\\scriptsize (service\_role / JWT)};
    \node[box, fill=colorAcento!30, draw=colorAcento!80!black] (rls) at (4.4,2.5) {Políticas RLS\\\scriptsize auth.uid() = dueño};
    \node[box, fill=c4Datos!25, draw=c4Datos!70!black] (db) at (4.4,0) {PostgreSQL\\\scriptsize esquema core};
    \node[box, fill=c4Componente!40] (rpc) at (4.4,-1.6) {RPC SECURITY\\DEFINER};

    \draw[rel] (cli) -- node[c4etiqueta, above]{SELECT directo} (rls);
    \draw[rel] (rls) -- node[c4etiqueta, right]{filas permitidas} (db);
    \draw[rel] (be) -- node[c4etiqueta, left]{jOOQ} (db);
    \draw[rel] (be) to[bend right=20] node[c4etiqueta, below]{invoca} (rpc);
    \draw[rel] (rpc) -- (db);
\end{tikzpicture}
\caption{Doble vía de acceso: el cliente pasa por RLS; el backend opera por jOOQ y RPCs.}
\label{fig:rls-flujo}
\end{figure}

\paragraph{Patrones de política.} Las $\sim$70 políticas RLS se reducen a unos pocos
patrones reutilizados, resumidos en la tabla~\ref{tab:patrones-rls}.

\vspace{0.2cm}
\noindent
\renewcommand{\arraystretch}{1.3}
\fontsize{8.2}{9.8}\selectfont\sffamily
\begin{tabularx}{\textwidth}{|>{\bfseries\color{colorPrimario}}l|>{\ttfamily\scriptsize}l|X|}
\hline
\rowcolor{headerTabla}
\textbf{Patrón} & \textbf{Expresión} & \textbf{Aplicación} \\
\hline
Dueño directo & profile\_id = auth.uid() & academic\_records, work\_experiences, cv\_documents, profile\_connections. \\
\hline
\rowcolor{grisTecnico}
Participante de chat & EXISTS(... chats ...) & chat\_messages: solo los dos participantes ven/insertan mensajes. \\
\hline
Lectura pública & visibility = 'Public' & projects, project\_media: lectura por \texttt{anon} solo de contenido público no borrado. \\
\hline
\rowcolor{grisTecnico}
Portafolio publicado & is\_published = true & portfolio\_settings: visible si está publicado. \\
\hline
Doble parte & company OR basic = uid() & profile\_likes, chats: visible para reclutador o candidato. \\
\hline
\rowcolor{grisTecnico}
Solo service\_role & true (service\_role) & verification\_codes, hard\_skills, soft\_skills: gestionados por el backend. \\
\hline
\end{tabularx}
\label{tab:patrones-rls}

A continuación, dos políticas reales que ejemplifican los patrones más sofisticados: la
visibilidad de mensajes restringida a los participantes del chat, y la nueva política
(v1.1.0) que permite a un candidato ver los \textit{likes} recibidos.

\begin{lstlisting}[language=SQL, caption={Políticas RLS representativas (esquema \texttt{core}).}]
-- Solo los participantes del chat ven sus mensajes
CREATE POLICY "participants can view messages" ON core.chat_messages FOR SELECT
USING (EXISTS (SELECT 1 FROM core.chats c
   WHERE c.chat_id = chat_messages.chat_id
     AND (c.company_profile_id = auth.uid() OR c.basic_profile_id = auth.uid())));

-- El candidato puede ver los likes que ha recibido (migración V15)
CREATE POLICY profile_likes_owner_select ON core.profile_likes FOR SELECT
USING (basic_profile_id = auth.uid());
\end{lstlisting}

\subsection{Gestión de roles de base de datos}
\label{subsec:roles}

El modelo de acceso de Supabase define tres roles principales sobre los que se evalúan
las políticas RLS. La tabla~\ref{tab:roles} describe cada uno y su uso en EthosHub.

\clearpage
\vspace{0.2cm}
\noindent
\renewcommand{\arraystretch}{1.3}
\fontsize{9}{10.8}\selectfont\sffamily
\begin{tabularx}{\textwidth}{|>{\ttfamily\bfseries\color{colorPrimario}}l|X|}
\hline
\rowcolor{headerTabla}
\sffamily\textbf{Rol} & \sffamily\textbf{Descripción y uso} \\
\hline
anon & Usuario \textbf{no autenticado}. Solo accede a lectura pública (portafolios publicados, proyectos \texttt{Public}). \\
\hline
\rowcolor{grisTecnico}
authenticated & Usuario \textbf{autenticado} (JWT válido). Las políticas comparan \texttt{auth.uid()} contra el dueño de la fila. \\
\hline
service\_role & Rol \textbf{privilegiado} del backend. Evade RLS (\texttt{USING true}) para operaciones de servidor; su clave nunca se expone al cliente. \\
\hline
\end{tabularx}
\label{tab:roles}

\begin{AlertaSeguridad}
La \comando{anon key} es \textbf{pública por diseño} y se expone en el cliente; su
seguridad reposa íntegramente en RLS. La \comando{service\_role key}, en cambio, evade
RLS y \textbf{nunca} debe llegar al navegador: vive solo en el backend como variable de
entorno. La separación estricta entre ambas claves es la piedra angular del modelo de
seguridad de datos de EthosHub.
\end{AlertaSeguridad}

\subsection{Síntesis de la cobertura de seguridad}
\label{subsec:sintesis_rls}

La tabla~\ref{tab:cobertura-rls} ofrece una visión consolidada de la cobertura de RLS
por tabla, evidenciando que ninguna tabla con datos de usuario queda desprotegida.

\vspace{0.2cm}
\noindent
\renewcommand{\arraystretch}{1.2}
\fontsize{8}{9.6}\selectfont\sffamily
\begin{tabularx}{\textwidth}{|>{\ttfamily\scriptsize\color{colorPrimario}}l|c|c|X|}
\hline
\rowcolor{headerTabla}
\sffamily\textbf{Tabla} & \sffamily\textbf{RLS} & \sffamily\textbf{Pol.} & \sffamily\textbf{Modelo de acceso} \\
\hline
profiles\_basic & Sí & 2 & Dueño (escritura) + lectura autenticada. \\ \hline
profiles\_company & Sí & 2 & Dueño + lectura autenticada. \\ \hline
\rowcolor{grisTecnico}
projects & Sí & 3 & Dueño + lectura pública + service\_role. \\ \hline
project\_media & Sí & 3 & Dueño vía proyecto + público + service\_role. \\ \hline
\rowcolor{grisTecnico}
work\_experiences & Sí & 4 & Dueño (CRUD completo). \\ \hline
academic\_records & Sí & 4 & Dueño (CRUD completo). \\ \hline
\rowcolor{grisTecnico}
hard\_skills / soft\_skills & Sí & 2 & Lectura autenticada + service\_role. \\ \hline
cv\_documents & Sí & 1 & Dueño (ALL). \\ \hline
\rowcolor{grisTecnico}
portfolio\_settings & Sí & 2 & Dueño + publicado. \\ \hline
profile\_connections & Sí & 2 & Dueño + service\_role. \\ \hline
\rowcolor{grisTecnico}
chats & Sí & 7 & Participantes (con borrado por parte). \\ \hline
chat\_messages & Sí & 7 & Participantes (ver/insertar/borrar propios). \\ \hline
\rowcolor{grisTecnico}
profile\_likes & Sí & 8 & Reclutador (CRUD) + candidato (ver recibidos). \\ \hline
recruiter\_candidate\_notes & Sí & 1 & Reclutador dueño. \\ \hline
\rowcolor{grisTecnico}
recruiter\_talent\_views & Sí & 4 & Reclutador dueño. \\ \hline
notifications & Sí & 2 & Dueño (ver/actualizar). \\ \hline
\rowcolor{grisTecnico}
verification\_codes & Sí & 1 & Solo service\_role (OTP blindado). \\ \hline
portfolio\_items & No & 0 & Acceso solo vía backend. \\ \hline
\rowcolor{grisTecnico}
talent\_score\_snapshots & No & 0 & Acceso solo vía funciones. \\ \hline
\end{tabularx}
\label{tab:cobertura-rls}

\clearpage
\begin{AvisoNota}
La tabla \comando{profile\_likes} concentra 8 políticas porque modela un recurso
compartido entre dos actores con derechos asimétricos: el reclutador gestiona el
\textit{like} (crear, mover de fase, eliminar), mientras que el candidato solo puede
\textbf{ver} los que ha recibido. Este matiz, introducido en la migración V15, es un
ejemplo de cómo RLS expresa reglas de negocio finas directamente en el motor.
\end{AvisoNota}
